\SourceLeaf{375}%
müsste \(DC=CL\) sein, im Widerspruch zum zuvor gefundenen Ergebnis. Wir glauben, dass wir bei diesem Schritt des Beweises den entscheidenden Punkt nach der Auffassung des Verfassers noch etwas klarer und bestimmter angedeutet haben als er selbst. Dadurch wird aber, wie uns scheint, auch seine Schwäche leichter erkennbar. Denn hier wird offenbar ganz willkürlich angenommen, dass bei allen gleichschenkligen Dreiecken mit dem Winkel \(A\) an der Spitze und größeren Winkeln an der Grundseite, wenn man die im Lehrsatz angegebene Konstruktion an ihnen ausführt, die Folge der abgeschnittenen Stücke im Hinblick darauf, ob sie gleich bleiben, größer oder kleiner werden, immer, unabhängig von der Größe der Seiten. notwendig dieselbe sein müsse. Diese Annahme kann unmöglich als selbstverständlich gelten. Da sich allein darauf der versuchte Beweis stützt, dass der dritte (wie auch der vierte und fünfte) Fall ausgeschlossen ist, und der ganze Artikel kein anderes Mittel bietet, den ausgelassenen Fall auszuschließen, glauben wir, damit das oben ausgesprochene Urteil hinreichend gerechtfertigt zu haben. Der gesamten übrigen einfallsreichen Durchführung in den folgenden Artikeln lassen wir aber gerne volle Gerechtigkeit widerfahren.

